\paragraph{How the remaining minors bound \(P+M\).}
For a fixed \((s,q)\), all the corresponding linear terms replace
columns at \(s+q\) by multiples of the same vector
\(\xi_s^{p^q}\) before the frame is applied. A determinant with two
such replacements vanishes. Therefore \(m_{sq}\le c_{s+q}\).
More generally the total number of replacements in a fixed set of
columns at a letter is at most the number of those columns in
the remaining minors, counted with multiplicity over the minors.
In a column set \(J\) the resulting bound over the remaining minors is
\[
  \#\{\text{replacements in }J\}\le\sum_{C}|C\cap J|.
\]
This follows by expanding the Hasse coefficient by columns;
each column of a determinant changes at most once. A minor with
a repeated original column already vanishes.

Orient every top--top term from the endpoint with the larger
exponent index \(u\) to the endpoint with the smaller index \(v\)
in \eqref{eq:tag-j2}. The indices have opposite parity since \(d\) is even.
Their difference \(u-v>0\) is congruent to the forward displacement
from the source \(s\) to the target \(t\).
In particular the directed displacement obeys
\[
  u-v\equiv t-s\pmod d,\qquad u-v\ge1,\qquad u-v\equiv1\pmod2 .
\]
 Let \(q\) be its least
positive residue modulo \(d\), so \(q\) is odd and belongs to
\(D_s\). Such a term uses at least \(p^q\) units at \(s\).
For fixed \((s,q)\), the number of these terms together with
\(m_{sq}\) is at most \(a_s/p^q\), by the available top degree.

When \(q\notin D_s\) only the linear terms enter, and their
number is bounded also by \(c_{s+q}\). Now
\[
 \sum_{q>0\ {\rm odd}}p^{-q}=p^{-1}\sum_{a\ge0}p^{-2a}=\frac{p}{p^2-1}=\gamma.
\]
Subtract the two counts just obtained in \eqref{eq:tag-j3}, for every
\((s,q)\) separately. This proves the contribution \(A\) in \eqref{eq:tag-j1}.
Every top--top term also uses one unit at an even top, as does
each linear term whose source is even. Together those uses are
at most \(\sum_{s\ {\rm even}}a_s\). The odd-source linear
terms cost at most
\(\sum_{s\in T\ {\rm odd}}\sum_{q>0\ {\rm odd}}\min(a_s/p^q,c_{s+q})\).
This proves the other index-zero estimate.

For \(i>0\), choose integral bases for the incident inclusion
whose reduction has the standard codimension-\(i\) plane.
Write \(J=\{1,\ldots,i\}\) for the scaled target coordinates.
On a minor of height \(b\) with column list \(C\), the exponent
contributed by the normalized incident operator is
\[
 \begin{cases}
    \min(b,i)-|C\cap J|,&s\text{ even},\\
    |C\cap J|,&s\text{ odd}.
 \end{cases}                                                       \tag{J4}\label{eq:tag-j4}
\]
The diagonal map scales the \(J\)-coordinates by \(p^{-1}\)
in an even polynomial factor and by \(p\) in an odd
contragredient factor. The missing even powers are supplied by
the target determinant twist and the normalization in \eqref{eq:tag-p5}:
\[
 d_i(j)+i\Delta_j
   =\sum_{s\ {\rm even}}\sum_{a\le i}(\lambda_{s,a}+w_j(s))
   =\sum_{C\ {\rm in\ even\ diagrams}}\min(b_C,i).
\]
The second equality counts the boxes in the first \(i\) rows,
column by column with multiplicity. The calculation is at the
neighbor in the chart, with the normalized theta section already
included. Multilinearity of minors preserves their spans under
integral changes of basis of these same lattices.

For every top--top term with column in \(J\), pay one from its
odd top degree, and for column outside \(J\), pay one from its
even top degree. The corresponding height-one factors of \eqref{eq:tag-j4}
charge exactly these degrees. For a linear term with even source
and column outside \(J\), pay from the same even top factor.
For one with even source and column in \(J\), pay from a
replacement in the odd residual minor at its destination. By
the column count already proved these payments fit in

\[
  \sum_C |C\cap J|
\]
 over the odd residual minors. No top
factor has been charged beyond its multiplicity in the
\(y\)-derivative. All unused exponents in \eqref{eq:tag-j4} are nonnegative
because the columns in a nonzero minor are distinct.
Thus \eqref{eq:tag-j4} accounts for \(P\) and the even-source terms in
\(M\). For each odd source the number of replacements is at most
the smaller of the available source degree and the destination
capacity. Summing the remaining charges gives
\[
 \sum_{s\in T\ {\rm odd}}\sum_{q>0\ {\rm odd}}m_{sq}
 \le \sum_{s\in T\ {\rm odd}}\sum_{q>0\ {\rm odd}}
                    \min(a_sp^{-q},c_{s+q}),
\]
as required for the bound \(G\) in \eqref{eq:tag-j1}.
The ultrametric inequality and the Gauss valuation in the frame
entries pass the estimate to every coefficient. Saturation of the
integral flag-minor span, proved in Section~\ref{sec:finite-valuations},
passes it to each functional integral on that span.

For a joint series the factor in the integral formal parameters
can involve all places. In any prescribed jet its finitely many
contributing monomials are products of local monomials with
integral coefficients. The bounds above hold for every such monomial.
On a second-embedding orbit a lower bound for all those monomials,
such as the cyclic-jet bound in Section~\ref{subsec:cycle-jets},
adds its own exponent. Summing the entries on the class
representatives takes one stabilizer denominator. This proves
\eqref{eq:tag-p7} from the claimed local chart shape, integral evaluations and the numerical comparisons 
\[
  g^0_{jv}\le A,
\]
 \(g^U_{jv}\le G\)
on the selected range.

In the large-characteristic range the
index-zero integrality proved above suffices for the base gain.
These implications use the full range of heights, top sets and
cycle lengths with the explicit digit loss in \eqref{eq:tag-j1}. Applying them
to the global period also requires the chosen split charts to
have this shape in the incident bases and the complementary
places to give the stated forced terms and denominator.
\label{cor:hecke-entry-capacities-1}

\end{proof}

\paragraph{Comparison with the numbers on the selected cycle.}
\begin{proposition}[Valuation gains and denominator exponents on the selected cycle]\label{prop:selected-cycle-numerical-comparisons}
Use the first-embedding cycles and endpoint weights of
Section~\ref{sec:host-edge-tests}. For \(H=\log_p(2a)\ge1\),
the prescribed gains satisfy \(g^0_{jv}\le A\) and \(g^U_{jv}\le G\).
Put \(C=w''+\eta+c_G+N_m\), let \(d\) be the cycle length and
\(T\) the set of tops. Write \(B_s^*=-\lambda_{s,k}\). For the denominator exponent in \eqref{eq:tag-p5} one has
\[
 P_{j,v}\le
 \begin{cases}
  \sum_{s\in T\ {\rm odd}}(w_j(s)-B_s^*)+K_P dC, & r=1,\ K_P\ge1,\\
  r\bigl(|T|(2b'+c_G)+dC\bigr), & r=r_G<k .
 \end{cases}
\] For \(p>2a\) the digit loss dominates both
would-be gains and the integral index-zero estimate applies.
All comparisons hold uniformly on the chosen family.
\end{proposition}

\begin{proof}
\proofstep{Identify the least odd displacements on the actual cycle.}

To avoid a conflict
with the number \(a_s\) of extracted columns, denote the first
weight at a top by \(A_s^*\) and the negative of its last weight
by \(B_s^*\). Use the cycle numbering of
Section~\ref{sec:host-edge-tests}, so that a forward step is
addition of the Frobenius displacement \(a_p\) in the Galois group of the auxiliary field.
Let \(d\) be the additive order of the displacement \(a_p\). Two positions give the same position on the Frobenius cycle precisely when
\[
  s a_p=t a_p\quad\Longleftrightarrow\quad d\mid(s-t),
\]
so the identification is injective on the cycle. Odd positions
are the odd top subset of the up-gain formula. In particular
each odd least distance in \(D_s\) gives an excluded odd
displacement in the sum over the cycle. An earlier visit to that
target would contradict leastness modulo the cycle length.

\proofstep{Compare the two capacity programs with the selected fills.}

In the sums over the selected cycles replace \(a_s\) at a top by
\(A_s^*-2c_G\), and the destination capacity \(c_{s+q}\)
by the filled weight \(w_j(s+q)+c_G\). The first substitution
decreases the argument since 
\[
  |\lambda_{s,2}|\le c_G;
\]
the second increases it since the retained first or second
base row has weight at most \(c_G\). These substitutions
can only decrease the positive-part summands of \(A\).
For \(G\), write the positive first sum as the sum of
\(a_s p^{-q}\) over all positive odd \(q\) and all tops.
After subtraction the summand for a positive odd \(q\) is
\[
  \begin{cases}
      (a_s p^{-q}-c_{s+q})_+, & s\text{ an odd top},\\
      a_s p^{-q}, & s\text{ an even top}.
  \end{cases}
\]
Decreasing \(a_s\) and increasing \(c_{s+q}\) can therefore only
decrease the up gain. This argument works even if a substituted first
argument is negative.

\proofstep{Bound the digit lengths by the prescribed error term.}

Write \(N_m=k+m(t+1)\), \(H=\log_p(2a)\), and suppose
\(H\ge1\). The size bounds for the auxiliary data give 
\[
  N_j\le N_m\le a,
\]
\(c_G\le a\), \(A_s^*\le2a\), \(w_j(s)\le4a\)
at a common late threshold. Here the rank grows polynomially in \(x\), while
\(a=x\exp(x^\epsilon)\) and the width satisfies \(c_G=o(b')\)
as in Section~\ref{sec:host-edge-tests}. The selected fill is at most
\(B_s^*+b'/x+\eta\) at a top, where \(B_s^*\le5b'/4\),
and at most \(w''+\eta\) elsewhere; \(w''\le b'\le a\)
and \(\eta\le a\). These are the common fill and offset sizes
of the first embedding.

Removing the first-row excess at a top
leaves all row lengths at most \(w_j(s)+c_G\), also on each
added block by nonnegativity of its earlier fill. Hence
\(a_s\le 3a\le(2a)^2\) and 
\[
  l_s\le5a^2\le(2a)^4.
\]
In particular
\[
              D\le 8dN_m^2H+9|T|H^2 .                            \tag{J5}\label{eq:tag-j5}
\]
Indeed the two digit lengths are bounded by \(3H,5H\). This
is covered by the prescribed loss for this cycle with both coefficients
at least nine. Subtracting that loss proves
\(g^0_{jv}\le A\), \(g^U_{jv}\le G\) with the selected fills.

\proofstep{Bound the denominator exponent in the line and plane ranges.}

Set \(C=w''+\eta+c_G+N_m\). In the line range \(r=1\)
the odd top contribution is exactly \(w_j(s)-B_s^*\);
each other odd letter costs at most \(w''+\eta+c_G\).
Since \(j(t+1)/2\le N_m\), this gives the upper bound

\[
  \sum_{s\in T\ {\rm odd}}(w_j(s)-B_s^*)+K_Pd C
\]
 with
\(K_P\ge1\). In the plane range \(r=r_G<k\), the tail of
length \(r\) omits the large first row. Every entry in the
tail is at most \(c_G\). At tops the filled weight is at most
\(2b'+\eta\) because \(x\ge2\); the other weights have the
non-top bound above. Write \(C_{\rm plane}\) for the resulting total tail cost. Summing \eqref{eq:tag-p5} with these bounds gives
\[
  C_{\rm plane}\le r\{|T|(2b'+c_G)+d C\}.
\]
 Both are the prices subtracted on
their respective ranges in \eqref{eq:tag-c4}. They are nonnegative, since
\(w_j(s)\ge B_s^*\) at each top.

\proofstep{Use the integral baseline above the size cutoff.}

When \(p>2a\), both valuation gains before subtracting the digit contribution are at most \(2a\gamma |T|\). Indeed the capacity is
nonnegative and the positive part of any substituted first
argument is at most \(2a\). For \(2a\ge e\), monotonicity of
\(y/\log y\) on \([e,\infty)\) gives

\[
  2a/p\le\log_p(2a).
\]
 Since \(\gamma\le2/p\) and
\(|T|\le d\), the first error term
\(K_1dN_m^2\log_p(2a)\), \(K_1\ge9\), already exceeds
both gains. At these primes take the index-zero integral bound
at every stage, absorbing every unused neighbor operator in the
integral input as in the proof of \eqref{eq:tag-p8}. The constants in \eqref{eq:tag-j5} and
these bounds for the denominator exponent are independent of the word, prime, stage,
number of tops and coefficient extension.
\label{cor:hecke-entry-capacities-2}
\end{proof}
